\documentclass[preprint,12pt,authoryear]{elsarticle}

\usepackage{amsmath,amssymb,amsthm}
\usepackage{booktabs}
\usepackage{array}
\usepackage{tabularx}
\usepackage{graphicx}
\graphicspath{{./}{figures/}}
\usepackage{enumitem}
\usepackage[colorlinks=true,urlcolor=blue,linkcolor=blue,citecolor=blue]{hyperref}
\emergencystretch=3em
\setlength{\tolerance}{2000}

\newtheorem{proposition}{Proposition}
\newtheorem{corollary}{Corollary}
\newcommand{\pos}[1]{\left[#1\right]_+}
\newcommand{\dconv}{\Delta b_{\mathrm{conv}}}
\newcommand{\Ups}{\Upsilon_c}
\newcommand{\Eceil}{E_{\mathrm{ceil}}}
\newcommand{\Acmin}{A_c^{\min}}
\newcommand{\tconv}{\tau_{\mathrm{conv}}}

\journal{Ecological Modelling}

\begin{document}

\begin{frontmatter}

\title{Aggregate biocapacity does not identify its composition: a two-land model of ecological-capital drawdown}

\author[1]{Amin Abaee\corref{cor1}}
\ead{amin_abaee@ut.ac.ir}
\cortext[cor1]{Corresponding author. ORCID: 0000-0002-0019-1842.}
\affiliation[1]{organization={Independent Researcher}, city={Tehran}, country={Iran}}

\begin{abstract}
\footnotesize
Biocapacity is usually reported as one total, although land types differ in yield and regeneration. A rise in the total can therefore conceal a decline in the land that supports future production. We use a two-land model with delays in regeneration and population response to study this possibility. When the supply available to people falls short of demand, the model converts ecological-capital land into fast-provisioning land. If fast land contributes more annual biocapacity per hectare than capital land, the transfer adds to the total even as capital-land area falls. The total alone cannot reveal that change in composition.

In a fixed-demand comparison—with capital quality at its maximum, no accumulated debt or yield boost, and restoration switched off—the reduced model has at most one balance point. The balance curve meets the minimum-area boundary at a baseline demand ceiling of about $1.14$\,gha without turning back (no fold), and the reduced system's slowest return rate stays nonzero (no critical slowing down). Above that ceiling, capital-land area approaches its minimum. When a shortfall disappears, conversion stops; without restoration, the lower area persists but does not keep declining. In a representative simulation, capital-land area falls $46.2\%$ and its modeled contribution to biocapacity falls $51.7\%$, while total model biocapacity rises $39\%$. These are conditional model results, not measurements of global physical stocks.

In the National Footprint and Biocapacity Accounts, Edition 2026 (1961--2023), total world biocapacity rises $24.9\%$, while biocapacity per person falls $52.7\%$ and the production footprint per person rises $16.7\%$. Cropland adds $24.06$ percentage points to the total change; Forest biocapacity subtracts $5.95$ points, and built-up land adds $7.44$ as an accounting class. No comparable global physical-stock series was assembled, so we report no physical-stock composite or decline finding. The SI specifies requirements for such a comparison but contains no physical-stock inventory. These results show why component values should accompany aggregate totals.
\end{abstract}

\begin{keyword}
\footnotesize
biocapacity \sep ecological footprint \sep land-use conversion \sep aggregation \sep identifiability \sep critical slowing down \sep delay differential equations \sep carrying capacity
\end{keyword}

\end{frontmatter}

\section{Introduction}\label{sec:intro}

Ecologists use \emph{carrying capacity} for the population that a resource base can sustain. It is difficult to estimate, and many models treat it as a fixed ceiling \citep{Cohen1995}. But when demand exceeds renewal, use can damage the resources that support future production and change that ceiling. This paper examines a related measurement question: can a total biocapacity number rise while part of the productive land base is shrinking?

The National Footprint and Biocapacity Accounts (NFA) report biocapacity as annual biological supply and the Ecological Footprint as annual demand on biologically productive area. Both are expressed in global hectares (gha). A gha is a hectare adjusted to the world's average biological productivity for a particular year; its annual basis is implicit \citep{GFNStandards2009,GFNGuidebook2019}. This adjustment combines area and productivity in one account value. A rise in the total cannot by itself show whether land became more productive, more land entered production, or the mix of land classes changed. Research on composite indicators calls this \emph{compensation}: a surplus in one component can offset a deficit in another \citep{Nardo2008,Becker2017,Fischer2022}. Related critiques address how the Ecological Footprint should be interpreted \citep{Blomqvist2013,GiampietroSaltelli2014,vdBerghGrazi2015,Galli2016}. Aggregate biocapacity remains a useful summary, but it should be read alongside its component values.

Human--environment models often summarize environmental change as a single stock or an overall balance. Examples include the World3 system \citep{Meadows1972,Meadows1974,Meadows2004}, models of Easter Island \citep{BranderTaylor1998,Anderies2000}, and coupled Earth-system frameworks \citep{Motesharrei2016,Donges2017,CalvinBondLamberty2018,HendersonLoreau2018}. Other studies debate whether human carrying capacity is fixed, variable, or a useful concept at all \citep{Cohen1995,vdBerghRietveld2004,Simon1996,Meyer1994,MeyerAusubel1999,Dasgupta2021}. Global demographic assessments often project population separately from environmental constraints \citep{UN2022}. Meanwhile, land-use studies document the conversion of forests, grasslands, and wetlands to cropland \citep{Haberl2007,Krausmann2013,Erb2017}. Together, this work establishes that land conversion matters. What remains unclear is what a single biocapacity total reveals when land classes differ in yield and regeneration.

Consider an orchard with two productive kinds of ground. Fast-provisioning ground—such as cropland and intensively managed systems—can yield a substantial annual harvest and be replenished from reserve land. Ecological-capital ground—such as forest, soil, wetland, and long-rotation systems—holds assets whose productive value and regenerative capacity are rebuilt more slowly. The harvest is an annual flow; the trees, soil, and other assets are stocks that support future harvests. When available annual provision does not meet demand, the model allows a shortfall to prompt conversion of ecological-capital ground into fast-provisioning ground. In orchard terms, meeting demand by cutting trees can raise the current harvest while drawing down the standing productive stock and may reduce later harvests. Because the aggregate harvest combines contributions from both productive books, this conversion can raise the total when fast ground contributes more per hectare, even as the capital book shrinks. Replanting can restore next year's crop without restoring the standing stock or all the conditions that support it. Neither orchard nor harvesters respond instantly: tree loss, soil recovery, and demographic adjustment unfold on different timescales. A stable or rising harvest can therefore coexist with a shrinking productive base. This analogy explains an information limit; it is not evidence that NFA accounts measure global physical stocks or a prediction of inevitable or abrupt ecosystem failure. In the two-book analogy, one record follows the annual harvest and the other follows the assets behind it.

We represent this distinction with two effective land classes. Fast-provisioning land—such as cropland and intensively managed systems—can provide a high annual yield. Ecological-capital land—such as forest, soil, wetland, and long-rotation systems—supports future production but regenerates differently and often more slowly. When a shortfall triggers conversion from capital land to fast land, total biocapacity can rise even as capital-land area $A_c$ falls. This is a conditional model result: NFA account categories do not map directly to $A_c$ or measure the physical stocks in the orchard analogy. Related material-ledger work calls the broader flow--stock masking a \emph{productivity illusion}; this paper studies its composition-specific form \citep{Abaee2026c}.\footnote{The labels conflict: PDF v50, README/checksum v49, and TeX header v48.}

A second analogy concerns what a service measure can tell us. An elevator rated for ten people may keep running for a time while carrying fourteen, but continued service does not show whether repeated overloading has reduced the cable's safety margin. The analogy separates observed output from the condition of the asset behind it. It does not predict an inevitable or sudden failure: in the reduced fixed-demand model, the balance curve reaches a specified area minimum without a fold.

We ask three questions. First, when can conversion raise total biocapacity while capital-land area $A_c$ falls, how long can this pattern last, and can the total reveal the area loss? Second, does the modeled drawdown behave like a tipping point—with a fold, hysteresis, or early-warning signals—or does it follow a different path? Third, what limits the drawdown, and what would allow the land area to recover?

The paper makes three contributions. First, it builds a two-land model in which area is conserved, capital-land quality changes separately, and shortfalls can trigger land conversion with a delay. Second, it gives the exact condition under which conversion raises total biocapacity and shows why the total alone cannot reveal how land is divided between the two classes. It also separates a unique balance point under fixed demand from the history-dependent path the model takes to reach it. Third, it decomposes component changes in NFA Edition 2026 and examines a separate crop-area/yield proxy, while stating clearly what these accounts can and cannot measure. The paper also gives five testable predictions. Parameters are representative, not fitted, so simulated paths illustrate conditional mechanisms rather than forecasts.